\bookchapter{บทที่ 57: เส้นทางขั้นสูง — สิทธิ์การเข้าถึง การแจ้งเตือน และสถาปัตยกรรมสมัยใหม่}

\chapterintro{ทำแอปออกมาได้เป็นแค่จุดเริ่มต้น บทนี้จะคุยเรื่องก้าวสำคัญจาก “ทำได้” ไปสู่ “ทำได้ดี”: จะขอสิทธิ์การเข้าถึงอย่างสวยงามได้อย่างไร จะส่งการแจ้งเตือนอย่างไรไม่ให้น่ารำคาญ จะยกสถาปัตยกรรมขึ้นอีกขั้นได้อย่างไร และ — ในปี 2026 จะจับมือกับ AI เขียนโค้ด Android อย่างไร ทุกข้อคือสิ่งที่คุณจะเจอทันทีในโปรเจกต์จริง}

\booksection{สิทธิ์ขณะรัน: เคาะประตูก่อนแล้วค่อยเข้า}

กล้อง, ตำแหน่งที่ตั้ง, สมุดรายชื่อ, การแจ้งเตือน … “เขตอ่อนไหว” เหล่านี้ Android กำหนดให้แอปต้องขอจากผู้ใช้ขณะรันแอป (ตั้งแต่ Android 6 เป็นต้นมา) กฎคือ: ประกาศก่อน (ใน Manifest) แล้วค่อยขอ (ในโค้ด) ถ้าผู้ใช้ปฏิเสธต้องรับมืออย่างสุภาพ ใน Compose ใช้ rememberPermissionState จากไลบรารีทางการ Accompanist แล้วสบายใจที่สุด:

\begin{lstlisting}[language=Kotlin]
// ประกาศใน Manifest ก่อน:
// <uses-permission android:name="android.permission.CAMERA" />

@Composable
fun CameraButton() {
    val permission = rememberPermissionState(
        permission = android.Manifest.permission.CAMERA
    )
    Button(onClick = {
        when {
            permission.status.isGranted -> openCamera()
            permission.status.shouldShowRationale -> showExplainDialog()
            else -> permission.launchPermissionRequest()
        }
    }) { Text("ถ่ายรูป") }
}
\end{lstlisting}

\begin{tcbwarn}{ข้อควรรู้: ถูกปฏิเสธแบบถาวรแล้วทำยังไง?}
เมื่อผู้ใช้ติ๊ก “ไม่ต้องถามอีก” แล้ว launchPermissionRequest จะไม่ขึ้นหน้าต่างถามอีก ตอนนั้นให้ชี้ทางให้ผู้ใช้ไปเปิดเองในการตั้งค่าระบบ (ใช้ Intent เด้งไปหน้าตั้งค่าของแอป) พร้อมอธิบายในหน้าจอว่า “สิทธิ์นี้เอาไว้ทำอะไร” ทำตัวนุ่มนวลขึ้นหน่อย ผู้ใช้จะยอมเปิดสิทธิ์มากขึ้นเยอะ — ข้อความตอนขอสิทธิ์ ควรค่าแก่การเขียนให้ดี

\end{tcbwarn}

\booksection{การแจ้งเตือน: บอกสิ่งที่มีประโยชน์ในเวลาที่ถูกต้อง}

การแจ้งเตือนคือช่องทางที่แอปเข้าถึงผู้ใช้ และเป็นฟีเจอร์ที่ผู้ใช้ “บล็อก” ง่ายที่สุด วินัย 3 ข้อในการส่งการแจ้งเตือน: ส่งแต่สิ่งที่มีประโยชน์ ให้ผู้ใช้ปิดได้ อย่าไปรบกวนกลางดึก การใช้งานพื้นฐาน:

\begin{lstlisting}[language=Kotlin]
// Android 13+ ต้องขอสิทธิ์ POST_NOTIFICATIONS ก่อน (ดูบทก่อนหน้า)

// 1. สร้างแชแนลแจ้งเตือน (Android 8+ ต้องมี)
val channel = NotificationChannel(
    "reminder", "เตือน", NotificationManager.IMPORTANCE_DEFAULT
)
val manager = getSystemService(NotificationManager::class.java)
manager.createNotificationChannel(channel)

// 2. ส่งการแจ้งเตือน
val notification = NotificationCompat.Builder(this, "reminder")
    .setSmallIcon(R.drawable.ic_notification)
    .setContentTitle("เตือนจดบัญชี")
    .setContentText("วันนี้ยังไม่ได้จดบัญชีนะ")
    .setPriority(NotificationCompat.PRIORITY_DEFAULT)
    .build()
if (ActivityCompat.checkSelfPermission(this,
        Manifest.permission.POST_NOTIFICATIONS) == PackageManager.PERMISSION_GRANTED
) {
    manager.notify(1, notification)
}
\end{lstlisting}

งานตามเวลา (เช่น “แจ้งเตือนให้จดบัญชีทุกวัน 20:00”) ให้ใช้ WorkManager — มันคือ “ผู้จัดการตารางงานเบื้องหลัง” ที่ทางการแนะนำ รับประกันว่างานจะถูกรันในจังหวะที่เหมาะสม ไม่ต้องไปกังวลเรื่องประหยัดแบตเอง งานแนวปลุกเตือนแบบนี้ใช้มันถูกทางแน่นอน

\booksection{สถาปัตยกรรมขั้นถัดไป: dependency injection และการแบ่งโมดูล}

พอโปรเจกต์ใหญ่ขึ้น ปัญหาหวาน ๆ สองเรื่องจะตามมาเยือน:

\begin{itemize}
  \item การสร้าง object รกขรุงขราง: ทุกหน้าจอต่างพากัน new ฐานข้อมูล new API ของตัวเอง — ใช้ Hilt (เฟรมเวิร์ก dependency injection) มาจัดการให้เป็นระบบ “ใครต้องการอะไร เอ่ยปากแล้วมีคนส่งให้ทันที”;
  \item โค้ดกองเป็นภูเขา: โปรเจกต์โมดูลเดียวยิ่งแก้ยิ่งช้า — แยกเป็นหลายโมดูล (app / core / feature-xxx) ให้แต่ละฟีเจอร์พัฒนาและทดสอบแยกจากกัน.
\end{itemize}

\begin{tcbinfo}{เคล็ดลับ: อย่าออกแบบเกินจำเป็น}
ทั้ง Hilt และการแบ่งโมดูลเป็นเครื่องมือดี แต่โปรเจกต์คนเดียวอย่ารีบใช้ แอปบันทึกรายรับ-รายจ่ายแบบโมดูลเดียว + สร้าง dependency ด้วยมือก็เอาอยู่สบาย รอโค้ดเกินหลายหมื่นบรรทัด คอมไพล์เริ่มช้า หรือต้องทำงานหลายคน ค่อยแนะนำเข้ามาตามความจำเป็น เครื่องมือมีไว้รับใช้เป้าหมาย ไม่ใช่ไว้โชว์ — นี่ก็คือหลัก “เรียบง่ายเหนือสิ่งอื่นใด” ในหนังสือเขียนโค้ดด้วย AI ของเรานี่เอง

\end{tcbinfo}

\booksection{วิธีเล่นใหม่ปี 2026: ใช้ AI ช่วยพัฒนา Android}

ปี 2026 AI ได้แทรกซึมเข้ากระบวนการพัฒนา Android อย่างลึกซึ้งแล้ว ถ้าอยากเร่งความเร็ว ชุดท่าต่อเนื่องนี้เข้าทางมาก:

\begin{center}
\begin{tabularx}{\linewidth}{|X|X|X|}
\hline
\textbf{ขั้นตอน} & \textbf{AI ช่วยอย่างไร} & \textbf{หน้าที่ของคุณ} \tabularnewline
\hline
เขียนโค้ด & Gemini ที่ฝังมาใน Android Studio / ปลั๊กอิน AI ของบุคคลที่สาม ช่วยเติมและสร้างโค้ด & อธิบายความต้องการให้ชัด และตรวจผลงาน \tabularnewline
ค้นคว้าข้อมูล & โยนเอกสารทางการกับข้อความ error ให้ AI ช่วยอ่านแปลความ & ข้อสรุปสำคัญต้องกลับไปเช็กกับเอกสารทางการ \tabularnewline
เขียนเทสต์ & ให้ AI สร้างเทสต์เคสตาม 3 ประเภท “ปกติ / ขอบเขต / ผิดปกติ” & ยืนยันว่าเทสต์เคสสอดคล้องกับความหมายทางธุรกิจ \tabularnewline
แก้บั๊ก & แปะ stack trace + โค้ดที่เกี่ยวข้อง ให้ AI ชี้จุดและแก้ & อ่านการแก้ไขให้เข้าใจ แล้วรันเทสต์ยืนยัน \tabularnewline
\hline
\end{tabularx}
\end{center}

\begin{center}\small\color{gray}4 สถานการณ์ของการใช้ AI ช่วยพัฒนา Android\end{center}

วิธีทำแบบละเอียด — จะสั่ง AI อย่างไร จะตรวจรับงานอย่างไร จะเลี่ยงหลุมพรางอย่างไร — มีครบ 11 บทในภาคที่ 7 “ลงมือเขียนโค้ดด้วย AI” (บทที่ 37-47) ของหนังสือเล่มนี้ แนะนำอย่างแรงให้อ่านคู่กัน: ไวยากรณ์ Android ฝากไว้กับภาคนี้ ส่วนศาสตร์การทำงานร่วมกับ AI ฝากไว้กับภาคที่ 7

\booksection{ปรับปรุงประสิทธิภาพ Compose กับเครื่องมือดีบัก}

เขียน Compose สนุกจริง แต่ “recompose เยอะเกินไป” ก็ทำให้กระตุกได้ สามเทคนิคที่ใช้ได้จริง:

\begin{itemize}
  \item ทำพารามิเตอร์ให้เสถียร: ใส่ @Immutable ให้ data class หรือใช้ derivedStateOf เพื่อกันไม่ให้สถานะที่ไม่เกี่ยวข้องเปลี่ยนแล้วจุด recompose ทั้งจอ;
  \item อย่าลืมใช้ key: items(list, key = \{ it.id \}) ทำให้เวลาเพิ่มหรือลบรายการ มีแค่รายการที่ควรขยับเท่านั้นที่ขยับ — ไม่ใส่ key แล้วลบไปหนึ่งชิ้น ทั้งลิสต์อาจวาดใหม่หมดเลย;
  \item โหมด Compose ของ Layout Inspector: ดูจำนวนครั้งที่แต่ละ Composable ถูก recompose แบบสด ๆ ใคร recompose ถี่เห็นกันชัด ๆ แล้วปรับปรุงตามจุดร้อน.
\end{itemize}

อีกอย่าง เปิดรายงานของ Compose compiler (Metrics) ด้วย จะได้รายงานว่า “คอมโพเนนต์ไหนยังไม่เสถียร” — บนเส้นทางขั้นสูง ข้อมูลเชื่อถือได้เสมอกว่าความรู้สึก

\begin{tcbpractice}{ลงมือทำ: เพิ่ม “ชุดสามชิ้น” ให้แอปบันทึกรายรับ-รายจ่าย}
(1) เพิ่ม “แจ้งเตือนให้จดบัญชีทุกวัน 20:00” ให้แอป (WorkManager + notification channel + สิทธิ์แจ้งเตือนของ Android 13); (2) เปลี่ยนการสร้างฐานข้อมูลและ API ให้เป็นการ inject ผ่าน Hilt; (3) ลองใช้เครื่องมือ AI เร่งทำฟีเจอร์ “กรองรายการรายเดือน” แล้วตรวจผลงานด้วยวิธีตรวจรับ 3 ชั้น ทำชุดสามชิ้นจบ โปรเจกต์ของคุณก็จะดูเป็น “ทีมงานมืออาชีพ” ขึ้นมาทันที

\end{tcbpractice}

\begin{tcbsummary}{สรุปบทนี้}
\begin{itemize}
  \item สิทธิ์อ่อนไหว: ประกาศใน Manifest + ขอขณะรัน + ถูกปฏิเสธแล้วชี้ทางอย่างสุภาพ
  \item วินัย 3 ข้อของการแจ้งเตือน: ส่งแต่ของมีประโยชน์ ปิดได้ ไม่รบกวน; ต้องมีทั้ง channel และสิทธิ์ ขาดอย่างใดอย่างหนึ่งไม่ได้
  \item งานตามเวลาใช้ WorkManager อย่าปั่นล้อขึ้นมาเอง
  \item Hilt ดูแล object, โมดูลดูแลโค้ด — โปรเจกต์เล็กคนเดียวอย่ารีบใช้
  \item AI ช่วยได้ 4 สถานการณ์: เขียนโค้ด ค้นข้อมูล เขียนเทสต์ แก้บั๊ก; วิธีทำอยู่ในภาคที่ 7 (บทที่ 37-47)
\end{itemize}

\end{tcbsummary}
